% Generated draft: requires mathematical and source review.
\documentclass[11pt]{article}
\usepackage[margin=25mm]{geometry}
\usepackage{fontspec}
\setmainfont{Latin Modern Roman}
\usepackage{amsmath,amssymb,mathtools,unicode-math}
\setmathfont{Latin Modern Math}
\usepackage{enumitem,needspace,xurl,hyperref}
\setlength{\parindent}{0pt}
\setlength{\parskip}{5pt}
\setlength{\emergencystretch}{3em}
\newcommand{\N}{\mathbb{N}}
\newcommand{\Z}{\mathbb{Z}}
\newcommand{\Q}{\mathbb{Q}}
\newcommand{\R}{\mathbb{R}}
\newcommand{\C}{\mathbb{C}}
\newcommand{\F}{\mathbb{F}}
\newcommand{\A}{\mathbb{A}}
\newcommand{\PP}{\mathbb{P}}
\newcommand{\E}{\mathbb{E}}
\newcommand{\eps}{\varepsilon}
\newcommand{\dd}{\,\mathrm{d}}
\begin{document}
\section*{UnsolvedMath ID: 30004274}
\textbf{Source problem code:} catalogue entry
\subsection{Definitions and mathematical statement}
The supplied problem concerns a three-dimensional equiregular sub-Riemannian structure locally modeled on the Heisenberg group and an operator obtained by deleting a point. Under the usual step-two interpretation, write this structure as \((M,H,g_H)\), where \(M\) is a smooth manifold, \(H\subset TM\) is a smooth distribution, and \(g_H\) is an inner product on \(H\), with
\[
\dim M=3,\qquad \operatorname{rank}H=2,\qquad H+[H,H]=TM.
\]
Here \([H,H]\) denotes the directions generated by brackets of local horizontal vector fields. The corresponding homogeneous dimension is
\[
Q=1\cdot\dim H+2\cdot\dim(TM/H)=2+2=4.
\]
The source does not specify whether being locally modeled on the Heisenberg group means local equivalence to its standard structure or only having its nilpotent approximation.

To specify the analytic question, one must additionally fix a positive smooth density \(\mu\) on \(M\) and a point \(p\in M\). Given these data, define the horizontal gradient by
\[
g_H(\nabla_H f,X)=\dd f(X)\qquad(X\in H),
\]
and define divergence by \(\mathcal L_X\mu=(\operatorname{div}_{\mu}X)\mu\). A standard interpretation of the pointed sub-Laplacian is the densely defined symmetric operator
\[
A_{p,\mu}f=-\operatorname{div}_{\mu}(\nabla_H f),\qquad
\mathcal D(A_{p,\mu})=C_c^\infty(M\setminus\{p\}),\qquad
A_{p,\mu}:\mathcal D(A_{p,\mu})\subset L^2(M,\mu)\longrightarrow L^2(M,\mu).
\]
This choice of density and initial domain is an explicit interpretation, not information supplied by the record. Essential self-adjointness means that the operator closure is self-adjoint, equivalently
\[
\overline{A_{p,\mu}}=A_{p,\mu}^{*},
\]
or that both deficiency spaces \(\ker(A_{p,\mu}^{*}-iI)\) and \(\ker(A_{p,\mu}^{*}+iI)\) vanish.

\textbf{Requested mathematical target.} Determine whether the topological dimension \(3\), horizontal dimension \(2\), or homogeneous dimension \(4\) governs essential self-adjointness of the pointed sub-Laplacian in the stated Heisenberg-type setting. The source supplies neither a definition of ``governs'' nor a quantified threshold criterion. Since all three dimensions are fixed in this setting, a self-adjointness determination alone does not specify a general dimensional criterion. Moreover, no density, completeness assumption, boundary conditions, or control of other ends or singularities is stated. Consequently, the record does not determine a unique global operator problem or a precise universal assertion without further qualifications.

The supplied literature assessment reports essential self-adjointness after deletion of a point in the three-dimensional genuine sub-Riemannian case, including the Heisenberg example, while reporting that no general criterion identifies the relevant dimension for all pointed sub-Laplacians. These are supplied dated curation claims; their theorem hypotheses and present validity have not been independently verified here.
\subsection{Short English statement}
The problem asks which notion of dimension controls uniqueness of the self-adjoint realization of a sub-Laplacian after removing a point from a three-dimensional Heisenberg-type space. The candidate dimensions are \(3\), \(2\), and \(4\). The supplied record labels the problem partially solved, but leaves the operator data and the meaning of a dimensional criterion unspecified. This is a statement-only draft requiring human mathematical review.
\subsection{Sources}
\begin{itemize}
\item Supplied local UnsolvedMath registry record, permanent ID \(30004274\), source problem code \texttt{catalogue entry}, title \emph{Dimensional Thresholds for Pointed Sub-Laplacian Self-Adjointness}; collection: Oberwolfach Reports Open Problems. The record gives proposed year \(2019\) and names no proposer. Its labels \texttt{partially\_solved} and \texttt{statement\_status: exact}, together with its assessment that the statement is self-contained and publication-ready, are supplied registry assertions, not independent mathematical verification.
\item \emph{Mini-Workshop: Self-adjoint Extensions in New Settings} (2020), Oberwolfach Reports, pp.\ 2911--2950, DOI \texttt{10.4171/owr/2019/47}: \url{https://doi.org/10.4171/owr/2019/47}. This is also the supplied literature source labeled ``original Oberwolfach report.'' The background attributes the problem collection to workshop open-problem discussions documented in Oberwolfach Reports; no exact page or contribution for this problem is supplied.
\item Adami--Boscain--Franceschi--Prandi, \emph{Point interactions for 3D sub-Laplacians}: \url{https://doi.org/10.1016/j.anihpc.2020.10.007}. Author attribution and title are reproduced from the supplied literature source.
\item Literature assessment checked at \texttt{2026-08-22}: supplied dated curation, not independently verified. It labels the problem partially solved, reports essential self-adjointness after point deletion in the three-dimensional genuine sub-Riemannian case including the Heisenberg example, and reports the absence of a general dimensional criterion for all pointed sub-Laplacians. The background describes this as dated literature triage derived from the supplied OWR curation and explicitly calls for independent verification of time-sensitive status and cited claims. No source visit is claimed here.
\end{itemize}
\end{document}
